\chapter{Tests, sécurité et déploiement}

\section*{Introduction}
Ce dernier chapitre d'implémentation présente la démarche de qualité logicielle adoptée pour la
plateforme PMS~: la stratégie de tests automatisés, la validation rigoureuse du contrôle d'accès
dynamique, la revue de sécurité au regard des risques de l'OWASP, et enfin la préparation au
déploiement.

\section{Stratégie de tests}

La priorité a été donnée aux \textbf{tests d'intégration} de niveau contrôleur, qui valident le
comportement réel de l'API --- routage, sérialisation, sécurité, règles métier et persistance ---
plutôt qu'à une multiplication de tests unitaires sur des composants isolés. Chaque test démarre un
contexte Spring complet avec une base \textbf{H2 en mémoire} et le profil \texttt{test}, et exerce
l'API via \texttt{MockMvc}.

Deux profils utilisateurs sont préparés par la fabrique \texttt{TestFixtures}~:
\begin{itemize}
  \item un \textbf{super-administrateur} doté de toutes les permissions, y compris la capacité
        d'accès portefeuille \texttt{VIEW\_ALL\_PROJECTS}~;
  \item un \textbf{lecteur} ne disposant que des permissions \texttt{VIEW\_*}, utilisé pour prouver
        que les opérations d'écriture sont bien refusées (\texttt{403}).
\end{itemize}

\begin{longtable}{|m{5.2cm}|m{1.8cm}|m{6.4cm}|}
\hline
\textbf{Suite de tests} & \textbf{Nombre} & \textbf{Portée principale} \\
\hline
\endhead
\endfoot
\endlastfoot
\hline
\texttt{AuthControllerTest} & 7 & Connexion, rafraîchissement, déconnexion, changement de mot de passe \\
\hline
\texttt{ProjectControllerTest} & 9 & CRUD, 401, \textbf{403 RBAC}, 404, 409 \\
\hline
\texttt{TeamControllerTest} & 6 & Affectation, doublon (409), retrait (204), historique \\
\hline
\texttt{WorkloadControllerTest} & 7 & Plan de charge, charge réelle, validation, idempotence \\
\hline
\texttt{BillingControllerTest} & 8 & Jalons, paiements, avenants \\
\hline
\texttt{KpiControllerTest} & 6 & KPI temps réel, instantanés, 201 + Location, 409, 404 \\
\hline
\texttt{GovernanceControllerTest} & 8 & Risques, machine à états des livrables, demandes de changement \\
\hline
\textbf{Total} & \textbf{51} & \textbf{51 / 51 réussis --- BUILD SUCCESS} \\
\hline
\caption{Couverture des tests d'intégration backend}
\label{tab:tests}
\end{longtable}

\noindent Les tests vérifient non seulement les chemins nominaux mais aussi, systématiquement, les
chemins d'erreur~: absence de jeton (\texttt{401}), permission insuffisante (\texttt{403}), entité
introuvable (\texttt{404}), doublon ou transition d'état invalide (\texttt{409}).

\section{Validation du contrôle d'accès dynamique}

Le contrôle d'accès étant le c\oe ur du projet, il a fait l'objet d'une validation dédiée en trois
temps.

\paragraph{Audit de la matrice d'autorisation.} La matrice rôle~$\rightarrow$~permission a été
\textbf{reconstruite à partir des règles métier} (et non du code existant), révélant plusieurs
incohérences de la version initiale~: un administrateur doté de droits opérationnels et financiers,
un développeur ayant accès aux indicateurs financiers (violation de BR-050), un directeur gérant la
facturation. La correction a réduit l'administrateur de 22 à 5 permissions et réaligné chaque rôle
sur sa responsabilité métier.

\paragraph{Une correction sans code.} Conformément à l'ADR-001, cette correction s'est faite par une
\textbf{unique migration de données} (\texttt{V12}) modifiant les associations rôle--permission,
\textbf{sans aucune ligne de code backend ou frontend}. C'est la démonstration la plus probante de la
valeur du RBAC dynamique~: un changement de politique d'autorisation est un changement de
configuration, pas de programme.

\paragraph{Vérification par rôle et par périmètre.} Deux campagnes de tests bout-en-bout ont confirmé
le résultat~: une grille \textbf{par rôle} (22~contrôles, chacun vérifiant un accès autorisé ou un
refus \texttt{403}) et une grille \textbf{de périmètre} (13~contrôles vérifiant que le directeur voit
tout, tandis que le chef de projet et le développeur sont restreints à leurs projets). Les deux
campagnes sont passées intégralement.

\section{Revue de sécurité (OWASP)}

La sécurité a été pensée dès la conception (ADR-001, ADR-010, ADR-017, ADR-021) puis revue au regard
des principaux risques de l'OWASP Top~10.

\begin{longtable}{|m{4.6cm}|m{9.4cm}|}
\hline
\textbf{Risque OWASP} & \textbf{Mesure dans PMS} \\
\hline
\endhead
\endfoot
\endlastfoot
\hline
A01 --- Contrôle d'accès défaillant & RBAC dynamique par permission \emph{et} périmètre de données
(permission~$\wedge$~scope)~; double application backend (\texttt{@PreAuthorize} + intercepteur de
périmètre) et frontend~; re-vérification de chaque identifiant fourni par le client. \\
\hline
A02 --- Défaillances cryptographiques & Mots de passe hachés en \textbf{BCrypt} (coût~12)~; jetons
JWT signés (HMAC-SHA512)~; aucun secret transporté dans les jetons. \\
\hline
A03 --- Injection & Accès aux données via JPA / requêtes paramétrées (aucune concaténation SQL)~;
schéma géré par Flyway~; validation des entrées par \emph{Bean Validation}. \\
\hline
A05 --- Mauvaise configuration & CORS restreint aux origines connues~; profils \texttt{dev}/\texttt{prod}
séparés~; \texttt{ddl-auto=validate} interdisant toute génération de schéma implicite. \\
\hline
A07 --- Défaillances d'authentification & Changement de mot de passe forcé au premier accès~;
révocation immédiate par \emph{token-version} (désactivation, changement de rôle/permission,
déconnexion)~; jeton d'accès de courte durée + rafraîchissement rotatif. \\
\hline
A09 --- Journalisation \& traçabilité & Champs d'audit (\texttt{createdBy/At}, \texttt{modifiedBy/At})
et suppression douce sur les entités majeures (ADR-009), préservant l'historique. \\
\hline
\caption{Couverture des principaux risques OWASP}
\label{tab:owasp}
\end{longtable}

\noindent Côté frontend, Angular échappe par défaut les interpolations, neutralisant les attaques
XSS (A03) les plus courantes. L'authentification étant portée par un en-tête \texttt{Authorization}
(et non par un cookie de session pour le jeton d'accès), la surface CSRF sur l'API est fortement
réduite~; le jeton de rafraîchissement, lui, est destiné à un cookie \texttt{HttpOnly~/~Secure}.

\section{Déploiement}

Le déploiement retenu est \textbf{natif}, sans conteneur (ADR-014), aligné sur l'environnement de la
machine de développement~:

\begin{itemize}
  \item \textbf{Base de données}~: service \textbf{PostgreSQL~17}~; le schéma est appliqué
        automatiquement au démarrage par les migrations \textbf{Flyway} versionnées
        (\texttt{V1}\dots\texttt{V15}).
  \item \textbf{Backend}~: application \textbf{Spring Boot} (Java~21) empaquetée en \texttt{jar}
        exécutable~; configuration externalisée par variables d'environnement et profils
        \texttt{dev}/\texttt{prod}.
  \item \textbf{Frontend}~: build de production Angular (\texttt{ng build}) générant des fichiers
        statiques servis par un serveur web, l'adresse de l'API étant paramétrée par fichier
        d'environnement.
\end{itemize}

\noindent La séparation stricte du schéma (Flyway, source de vérité) et des entités
(\texttt{ddl-auto=validate}) garantit qu'aucune dérive ne peut survenir entre le code et la base~:
un écart fait échouer le démarrage plutôt que de corrompre silencieusement les données.

\section*{Conclusion}
La plateforme PMS est validée par 51~tests d'intégration et deux campagnes de vérification du
contrôle d'accès, sécurisée selon les bonnes pratiques de l'OWASP, et prête à un déploiement natif
reproductible. La qualité n'est pas un ajout final mais une propriété construite phase après phase~:
permissions dynamiques, périmètre de données, suppression douce, migrations versionnées et tests
systématiques forment un ensemble cohérent. La conclusion générale dresse le bilan du travail
accompli et ouvre sur les perspectives d'évolution.
